%=====================================================================
\section{Lecture 4: Real-time diagrams for the reduced density matrix}
\label{sec:lec4}

We now combine the Dyson series \eqref{eq:dyson-super}, the vertex \eqref{eq:LIcq} and the contractions \eqref{eq:contractions} into a diagrammatic expansion of the dynamical map $\Pi(t)$, and evaluate it for the qubit.
The second-order result is the Born kernel that every master equation uses; the fourth-order structure introduces the notion of irreducibility on which \cref{sec:lec5} rests.

\subsection{Feynman rules}
\label{sec:feynman-rules}

Insert \cref{eq:LIcq} into \cref{eq:dyson-super} and take $\Tr_B$.
Every vertex is either $S^\qu B^\cl$ or $S^\cl B^\qu$, so the bath index at a vertex is the \emph{opposite} of the system index.
Apply Wick's theorem to the bath string and translate \cref{eq:LO} to the system side:

\begin{keybox}[System-side Feynman rules for $\Pi(t)$]
A term of order $2n$ is a time-ordered string of $2n$ vertices at $t>t_1>\dots>t_{2n}>0$, with a pairing of the vertices into $n$ lines.
Reading from right (earliest) to left:
\begin{itemize}[itemsep=1pt]
\item each vertex contributes $-i\,\tilde S^\qu(t_i)$ or $-i\,\tilde S^\cl(t_i)$ (interaction picture; in the Schr\"odinger picture, replace $\tilde S^a(t_i)$ by $S^a$ and insert free propagators $\ee^{\Lc_S(t_i-t_{i+1})}$ between consecutive vertices and $\ee^{\Lc_St_{2n}}$ before the first, $\ee^{\Lc_S(t-t_1)}$ after the last);
\item a \emph{noise line} $S(t_i-t_j)$ connects $S^\qu(t_i)$ to $S^\qu(t_j)$ ($\qu$--$\qu$ on the system side is $\cl$--$\cl$ on the bath side);
\item a \emph{response line} $-i\chi(t_i-t_j)$ connects the \emph{later} $S^\qu(t_i)$ to the \emph{earlier} $S^\cl(t_j)$;
\item there is no line with both ends $\cl$, and no line whose later end is $\cl$.
\end{itemize}
Consequently the latest vertex of every diagram is $S^\qu$, which preserves $\Tr_S$ (trace preservation of $\Pi$), and the earliest vertex of every line can be either.
\end{keybox}

The last rule has a direct physical reading.
A $\cl$ vertex is where the bath \emph{acts} on the system (a force, through $\chi$); a $\qu$ vertex is where the system \emph{acts} on the bath or where noise enters.
Causality says the bath can only respond to earlier system history, hence the response line points from an earlier $\cl$ to a later $\qu$.
For fermionic baths the same rules hold with a sign for each crossing of lines; for non-Gaussian baths, $2n$-point cumulants join more than two vertices~\cite{Schoeller2009}.

\subsection{Second order: the Born kernel}

\begin{exbox}[Example 4.1: the second-order term]
From \cref{eq:rho2super} with $\Lc_I=-i(S^\qu B^\cl+S^\cl B^\qu)$,
\begin{align}
  \Tr_B\big[\tilde\Lc_I(t_1)\tilde\Lc_I(t_2)\,\rho_S\!\otimes\!\rhoB\big]
  =-\Big[&\tilde S^\qu_1\tilde S^\qu_2\,\underbrace{\Tr_B[B^\cl_1B^\cl_2\rhoB]}_{S(t_1-t_2)}
   +\tilde S^\qu_1\tilde S^\cl_2\,\underbrace{\Tr_B[B^\cl_1B^\qu_2\rhoB]}_{-i\chi(t_1-t_2)}\notag\\
   &+\tilde S^\cl_1\tilde S^\qu_2\,\underbrace{\Tr_B[B^\qu_1B^\cl_2\rhoB]}_{0}
   +\tilde S^\cl_1\tilde S^\cl_2\,\underbrace{\Tr_B[B^\qu_1B^\qu_2\rhoB]}_{0}\Big]\rho_S ,
\end{align}
with the shorthand $\tilde S^{a}_i=\tilde S^{a}(t_i)$, $B^{a}_i=\tilde B^{a}(t_i)$.
Only two of the four index combinations survive, exactly as the rules say.
In the Schr\"odinger picture the surviving bracket, with $\tau=t_1-t_2$, defines the \emph{second-order kernel}
\begin{equation}
  \boxed{\;K^{(2)}(\tau)=-S^\qu\,\ee^{\Lc_S\tau}\Big[S(\tau)\,S^\qu-i\,\chi(\tau)\,S^\cl\Big],\;}
  \label{eq:K2}
\end{equation}
and the second-order map is $\Pi^{(2)}(t)=\int_0^t\dd t_1\int_0^{t_1}\dd t_2\,\ee^{\Lc_S(t-t_1)}K^{(2)}(t_1-t_2)\,\ee^{\Lc_St_2}$.
Expanding $S^\cl$ and $S^\qu$, \cref{eq:K2} is the textbook Born kernel $K^{(2)}(\tau)X=-\big[S,\ee^{\Lc_S\tau}\big(C(\tau)SX-C(-\tau)XS\big)\big]$~\cite[Sec.~3.3.1]{BreuerPetruccione}, but the rotated form keeps the two bath inputs apart, and that is what we want.
\end{exbox}

\begin{exbox}[Example 4.2: $K^{(2)}$ for the qubit, in Bloch form]
Apply \cref{eq:K2} to the coefficient basis $(\Id,\sx,\sy,\sz)$ using the matrices \eqref{eq:SclSqu} and the rotation \eqref{eq:LS-bloch}.
Write $c=\cos\theta$, $s=\sin\theta$, $S_\tau=S(\tau)$, $\chi_\tau=\chi(\tau)$, and $\cos,\sin$ for $\cos\omega_0\tau,\sin\omega_0\tau$.
Going through the four input vectors (the computation is written out in the solution to Exercise~\ref{ex:4-2}) gives
\begin{equation}
  K^{(2)}(\tau)=\scalebox{0.82}{$\displaystyle
  \begin{pmatrix}
  0 & 0 & 0 & 0\\[2pt]
  2cs\,\chi_\tau\sin & -4c^2S_\tau\cos & 4c^2S_\tau\sin & 4cs\,S_\tau\cos\\[2pt]
  2cs\,\chi_\tau(1\!-\!\cos) & -4c^2S_\tau\sin & -4S_\tau(c^2\cos+s^2) & 4cs\,S_\tau\sin\\[2pt]
  -2s^2\chi_\tau\sin & 4cs\,S_\tau\cos & -4cs\,S_\tau\sin & -4s^2S_\tau\cos
  \end{pmatrix}$},
  \label{eq:K2-bloch}
\end{equation}
Everything one needs to know about second-order qubit relaxation is in this matrix:
\begin{itemize}[itemsep=1pt]
\item The first row vanishes: trace preservation (last vertex is $S^\qu$).
\item The first \emph{column} is the only place $\chi$ appears.
  It multiplies the input $\Id$, i.e.\ it is the inhomogeneous term that drives the Bloch vector toward a nonzero steady state.
  Noise alone ($S$) would relax the qubit to the maximally mixed state; it is the \emph{response} of the bath that knows about the sign of $\omega_0$ and produces the Boltzmann asymmetry between emission and absorption.
\item Populations ($z$) relax only through $s^2=\sin^2\theta$: the transverse part of the coupling, sampled at the qubit frequency ($S_\tau\cos\omega_0\tau$).
\item The $c^2=\cos^2\theta$ terms in the coherence block are sampled at zero frequency ($S_\tau$ alone, after the rotating-frame average): pure dephasing.
  The $s^2$ term in $K_{yy}$ is the ``$T_1$ contributes to $T_2$'' piece.
\item The $cs$ cross terms couple populations to coherences.
  In the rotating frame they oscillate at $\pm\omega_0$ and average out when $\omega_0\gg$ rates (secular approximation, \cref{sec:lec5}); at short times they are real and are the reason Exercise~\ref{ex:2-4} has no simple answer.
\end{itemize}
The matrix is real, as Hermiticity preservation requires in this basis.
\end{exbox}

\subsection{Pure dephasing to all orders; reducible and irreducible diagrams}

\begin{exbox}[Example 4.3: $\theta=0$ and the first irreducible diagram]
Take $\theta=0$, so $S=\sz$ and $S^{\cl/\qu}$ commute with $\Lc_S$.
On the coherence $X=\ket0\bra1$ we have $S^\qu X=[\sz,X]=2X$ and $S^\cl X=\tfrac12\{\sz,X\}=0$.
Every response line ends on an $S^\cl$ and therefore vanishes on $X$.
What remains is a sum over pairings of noise lines, with every vertex equal to the number $-2i$:
\begin{align}
  \tilde\rho_{01}(t)=\rho_{01}(0)\Big[1
  &-4\!\int_{t_1>t_2}\!S_{12}
  +16\!\int_{t_1>t_2>t_3>t_4}\!\big[\underbrace{S_{12}S_{34}}_{\text{reducible}}
    +\underbrace{S_{13}S_{24}}_{\text{crossed}}
    +\underbrace{S_{14}S_{23}}_{\text{nested}}\big]-\dots\Big],
  \label{eq:dephasing-4th}
\end{align}
with $S_{ij}\equiv S(t_i-t_j)$, where the sign at order $2n$ is $(-2i)^{2n}=(-4)^n$.
Because all vertices commute, the sum over pairings is $\frac1{n!}$ times the $n$-th power of the single contraction, and the series exponentiates:
\begin{equation}
  \rho_{01}(t)=\rho_{01}(0)\,\ee^{-i\omega_0t}\,\ee^{-\Gamma(t)},\qquad
  \Gamma(t)=4\!\int_0^t\!\dd t_1\!\int_0^{t_1}\!\dd t_2\,S(t_1-t_2)
  =\frac4\pi\int_0^\infty\!\dd\omega\,\frac{S(\omega)}{\omega^2}\big(1-\cos\omega t\big).
  \label{eq:dephasing-exact}
\end{equation}
(The factor 4 is the ``our $S=\sz$ versus Leggett's $\sz/2$'' factor of \cref{sec:lec2}.)
This is the exact pure-dephasing result~\cite[Sec.~4.2]{BreuerPetruccione}; we will rederive it from the influence functional in \cref{sec:lec7} and use it in \cref{sec:lec10}.

Now look at the three fourth-order pairings.
The first, $S_{12}S_{34}$, can be cut between $t_2$ and $t_3$ without cutting a line: the bath has returned to $\rhoB$ there, and the diagram is a product of two second-order pieces.
It is \emph{reducible}.
The crossed and nested pairings cannot be cut: they are \emph{irreducible}, and they are the fourth-order contributions to the memory kernel of \cref{sec:lec5}.
Summing all orders, the reducible structure is what makes the exponential; the irreducible pieces are what the exponent is made of.
The general statement is \cref{sec:nz-irreducible}.
\end{exbox}

\begin{figure}[t]
\centering
\begin{tikzpicture}[>=Stealth,scale=1.0]
  \tikzset{v/.style={circle,fill=black,inner sep=1.3pt}}
  \draw[->] (0,0)--(5.2,0);
  \foreach \x in {0.6,1.6,3.0,4.2} \node[v] at (\x,0) {};
  \draw[thick,ForestGreen] (0.6,0) .. controls (0.8,0.8) and (1.4,0.8) .. (1.6,0);
  \draw[thick,ForestGreen] (3.0,0) .. controls (3.2,0.8) and (4.0,0.8) .. (4.2,0);
  \draw[dashed,red] (2.3,-0.3)--(2.3,0.9);
  \node[below,font=\small] at (2.6,-0.35) {reducible: $S_{12}S_{34}$};
  \begin{scope}[xshift=5.6cm]
  \draw[->] (0,0)--(5.2,0);
  \foreach \x in {0.6,1.9,3.0,4.2} \node[v] at (\x,0) {};
  \draw[thick,ForestGreen] (0.6,0) .. controls (1.0,1.1) and (2.6,1.1) .. (3.0,0);
  \draw[thick,ForestGreen] (1.9,0) .. controls (2.3,1.1) and (3.8,1.1) .. (4.2,0);
  \node[below,font=\small] at (2.6,-0.35) {crossed: $S_{13}S_{24}$};
  \end{scope}
  \begin{scope}[xshift=11.2cm]
  \draw[->] (0,0)--(5.2,0);
  \foreach \x in {0.6,1.6,3.0,4.2} \node[v] at (\x,0) {};
  \draw[thick,ForestGreen] (0.6,0) .. controls (1.0,1.3) and (3.8,1.3) .. (4.2,0);
  \draw[thick,ForestGreen] (1.6,0) .. controls (1.9,0.7) and (2.7,0.7) .. (3.0,0);
  \node[below,font=\small] at (2.6,-0.35) {nested: $S_{14}S_{23}$};
  \end{scope}
\end{tikzpicture}
\caption{The three fourth-order diagrams.
Time runs to the right; vertices are system superoperators; lines are bath contractions ($S$ or $\chi$).
Only the first can be cut (red) without cutting a line.}
\label{fig:fourth-order}
\end{figure}

\subsection{What the diagrams are good for}

The series \eqref{eq:dyson-super} converges for any finite $t$ but is useless as it stands: at order $2n$ there are $(2n-1)!!$ pairings and the terms grow like $(g^2t^2)^n$.
Every practical method is a \emph{resummation}:
\begin{itemize}[itemsep=1pt]
\item \textbf{Dyson resummation} (\cref{sec:lec5,sec:lec6}): sum the reducible structure exactly and approximate the irreducible kernel.
  This is Nakajima--Zwanzig, and TCL is the same thing organized differently.
\item \textbf{Exponentiation} (\cref{sec:lec7}): when the vertices commute (or are replaced by c-number paths), the pairings exponentiate into an influence functional.
  The noninteracting-blip approximation is a controlled truncation of that exponent.
\item \textbf{Self-consistent resummation} (\cref{sec:lec9}): replace bare propagators by dressed ones inside the kernel (Kadanoff--Baym, GKBA, DMFT-like schemes).
\end{itemize}
In all three, the only input from the bath is the pair $(S,\chi)$, and the only thing that distinguishes the methods is which subset of diagrams is kept and how it is organized.

\subsection{Exercises}

\begin{exercise}\label{ex:4-1}
From the Feynman rules, show that $\Tr_S K^{(2)}(\tau)X=0$ for every $X$ and interpret it.
Show also that $K^{(2)}(\tau)\rho$ is Hermitian for Hermitian $\rho$ (use $(S^\cl)^\dagger=S^\cl$, $(S^\qu)^\dagger=-S^\qu$ and the reality of $S(\tau),\chi(\tau)$).
\end{exercise}

\begin{exercise}\label{ex:4-2}
Derive \cref{eq:K2-bloch}.
(Apply \cref{eq:K2} to each of $\Id,\sx,\sy,\sz$ in turn; the solution writes out the $\sz$ and $\Id$ columns.)
Check the $\theta=0$ and $\theta=\pi/2$ limits against your expectations for pure dephasing and for a damped qubit.
\end{exercise}

\begin{exercise}\label{ex:4-3}
List the $15$ pairings at sixth order and classify them as reducible (cuttable) or irreducible.
How many are irreducible?
(Answer: 10 of the 15 are irreducible if one counts only cuts between consecutive vertices; the reducible ones are products of lower-order irreducible blocks.)
\end{exercise}

\begin{exercise}\label{ex:4-4}
Evaluate $\Gamma(t)$ in \cref{eq:dephasing-exact} for the Lorentzian bath at $T=0$ and show that $\Gamma(t)\simeq 2g^2t^2$ for $t\ll1/\Omega$ (Gaussian decay) and $\Gamma(t)\simeq2S(\omega{=}0)\,t+\text{const}$ for $t\gg1/\gamma$ (exponential decay).
Which value of $S(\omega{=}0)$ enters?
\end{exercise}

\clearpage
